\documentclass[a4paper,12pt]{ctexart}
\usepackage{xcolor}
\usepackage{graphicx}
\usepackage[top=2cm, bottom=2cm, left=2cm, right=2cm]{geometry}
\usepackage{array}
\usepackage{multirow}
\usepackage{hyperref}
\usepackage{framed}
\usepackage{enumitem}
\hypersetup{colorlinks=true,linkcolor=blue!70!black}
\linespread{1.5}

\newcommand{\ruby}[2]{#1\textsuperscript{#2}}
\newcommand{\xuehai}[2]{%
  \vspace{0.3cm}
  \noindent
  \fcolorbox{blue!50!black}{blue!5}{%
    \begin{minipage}{0.92\textwidth}
    \textbf{\textcolor{blue!60!black}{核心概念} · #1} \\[0.15cm]
    {\small #2}
    \end{minipage}%
  }
  \vspace{0.3cm}
}

\title{\textcolor{blue!70!black}{\Huge\textbf{欧洲古代史}}}
\author{}
\date{}

\begin{document}
\maketitle
\thispagestyle{empty}

\tableofcontents
\newpage

\section{第一课\quad 查理曼与欧洲的"分裂基因"}

\subsection{法兰克人的崛起}

公元5世纪中叶，西罗马帝国正在缓慢崩塌。在不列颠，罗马军团已于410年撤离，盎格鲁-撒克逊人乘船渡海而来；在高卢，法兰克人从莱茵河东岸向西推进；在西班牙，西哥特人建立了自己的王国；在北非，汪达尔人攻陷了迦太基，以该城为基地建立了一支地中海舰队，一度截断了罗马城从非洲进口粮食的海上粮食通道。476年，日耳曼雇佣军首领奥多亚塞废黜了最后一位西罗马皇帝罗慕路斯·奥古斯都路斯，将皇权标志送往君士坦丁堡。西罗马帝国——作为一个完整的政治实体——在地中海西方消失了。

但罗马留下了三样东西，盖过了帝国本身的废铁。第一，拉丁语：它维持了作为教会和学术的共同书写媒介的功能。任何一个受过拉丁文教育的教士，不论出生在今天是意大利、法兰西还是英格兰的地方，都能用同一套文字系统进行阅读、书写和辩论。第二，罗马法：虽然后来的蛮族王国各自颁布了习惯法汇编，但罗马法的术语和论证方式在各大学的法学院中被持续研究，并最终成为欧洲大陆大多数国家法律体系的基础。第三，基督教会：教区网络覆盖了帝国全境及其周边，修道院作为知识的保存者，是罗马帝国与中世纪之间最坚固的文化桥梁。

在所有日耳曼王国中，政治寿命最长的是法兰克人。法兰克人在5世纪末由克洛维一世统一并皈依了正统拉丁基督教（而非当时许多日耳曼部落信奉的阿利乌派异端基督教）。这意味着法兰克国王从一开始就和罗马教会站在了同一宗教战线上——这与当时控制着意大利、高卢南部和西班牙的东哥特王国和西哥特王国形成了鲜明对比（它们信奉阿利乌派，与教皇的关系长期处于紧张状态）。这一宗教一致性为法兰克王国与教皇之间的长期合作奠定了基础。

公元732年，法兰克宫相查理·马特在图尔-普瓦捷战役中击退了从西班牙北上的阿拉伯军队。这场战役被后世的欧洲编年史高度神话化，但它确实标志着一个转折点：伊斯兰势力在伊比利亚半岛的扩张被遏制在了比利牛斯山脉一线。查理·马特在这场战役中成功的关键在于他建立了一支以重装骑兵为核心的机动部队——骑士的马匹、铠甲和武器费用来自被没收的教会土地，这种"用土地换军事服役"的做法后来发展成了欧洲封建制度的经典模式。

弗兰克王国的版图在查理·马特的孙子查理曼统治期间达到了古代中世纪欧洲任何政治实体从未企及的最高峰。

\subsection{公元800年：一顶改变了欧洲的皇冠}

查理曼在768年继承法兰克王位时，他的王国已经涵盖了今天的法国、低地国家和德意志西部。他在位四十六年间发动了五十三次军事远征——最艰难的是对萨克森人的十八年战争（772—804），他最终用武力和强制改信将日耳曼腹地的最后一支独立部族纳入了基督教世界。萨克森战争也是查理曼统治中最具争议的篇章：据法兰克编年史记载，在一次大规模的萨克森人反抗后，查理曼在凡尔登处决了四千五百名萨克森俘虏。这是一个无法从法兰克史料之外获得独立验证的数字，但它清楚地传递了一个信息：查理曼的"基督教化"是以铁和血推进的。

当弗兰克王国的版图扩大到足以与昔日罗马帝国在西方的疆界相提并论时，查理曼的宫廷学者开始构想一个更宏大的合法性叙事：如果帝国恢复到查理曼手中，那么他应当被授予相应的称号。教皇利奥三世迫切需要查理曼的保护——他于799年在罗马被政敌袭击并囚禁，逃出后越过了阿尔卑斯山前往帕德博恩向查理曼求助。查理曼派兵护送他返回罗马并恢复了教宗职位。

作为回报，利奥三世在800年至诞节弥撒中，当查理曼跪在圣彼得大教堂的祭坛前时，将一顶皇冠戴在了他头上，在场的罗马教众齐声欢呼"查理·奥古斯特，罗马人的伟大皇帝，万寿无疆！"这一加冕从教皇的角度看是在确立教廷有权授予帝国头衔的先例；从查理曼的角度看，则是他收复了自476年以来中断的"罗马皇帝"的法统。两人的意图并不一致，但这一事件的后果是铁定的：西罗马帝国灭亡三百多年后，西欧再一次出现了一个自称继承罗马法统的帝国。

查理曼帝国全盛时期疆域东起易北河、西至大西洋、南抵意大利中部和西班牙边境、北达北海。首都设在亚琛（今德国西部），在那里查理曼召集学者重建宫廷学校，推广书写和拉丁文教育——史称"加洛林文艺复兴"。英格兰学者阿尔昆主持的宫廷学校系统地抄写和传播了大量罗马和早期基督教文献，这些文本的保存直接影响了后来的欧洲学术复兴。查理曼还任命"巡按使"成对巡察各个伯爵领地，监督地方行政和司法。加洛林时期的行政文书保留了许多罗马行政技术，包括使用标准化的公文格式和保存存档副本。

\subsection{凡尔登与继承的制度化难题}

查理曼于814年去世。他的儿子虔诚者路易继承了整个帝国，但路易的统治长期被几个儿子之间争夺继承权的内战所困扰。843年，路易的三个儿子——洛泰尔、日耳曼人路德维希和秃头查理——在凡尔登签署了瓜分帝国的条约。长子洛泰尔得到了从北海经阿尔卑斯山到意大利的狭长地带（"中法兰克"），日耳曼人路德维希得到了莱茵河以东的土地（东法兰克，德意志的雏形），秃头查理得到了莱茵河以西的土地（西法兰克，法兰西的雏形）。

\xuehai{继承制度的地缘政治后果}{
  法兰克人实行的继承规则是诸子平分，而非中国古代的嫡长子继承制。这一制度来自日耳曼部落的传统：父亲的遗产是所有儿子共有的，每个儿子都应按份继承。帝国被视为国王的个人财产，同样适用这一规则。凡尔登条约不是三代人之间情感破裂的偶然产物，而是这一继承制度的逻辑必然。

  数学上的后果是累积性的：一个王朝分一次，再分一次，再分一次——权力碎片化了。中国自周代就确立了嫡长子继承制，王国不可分割。秦以后更推广了中央任命的郡县制，在制度上杜绝了分裂的合法性。欧洲的诸子平分制度在一千年的时间尺度内，产生了截然不同的政治版图——分裂不是意外，是规则的必然产物。
}

中法兰克的路易二世死后无嗣，中法兰克在870年的墨尔森条约中再次被东西法兰克瓜分。从此，作为独立政治实体的"中法兰克"从地图上消失，它的核心地带——阿尔萨斯、洛林和低地国家——从此成为法德之间长达一千年的争夺对象。

\subsection{帝国的碎片化与封建制度的形成}

9至10世纪，查理曼的后代在相互攻伐中消耗殆尽。与此同时，来自三个方向的入侵潮席卷了西法兰克和东法兰克的疆域：南边的阿拉伯海盗袭击地中海沿岸，东边的马扎尔人（匈牙利人的前身）骑兵深入德意志和意大利，最致命的是来自北方的维京人的劫掠。

维京人的长船吃水极浅，可以沿河流深入内陆数百公里。845年，维京人的舰队沿塞纳河而上，攻入巴黎，西法兰克国王秃头查理支付了七千磅白银的赎金才让他们离开。在无力组织全国性防御的情况下，各地领主开始在各自的领地修建城堡、组建私人武装、供养自己的骑士。加洛林王朝的中央权力在一次次入侵面前持续收缩，权力碎片化到地方层面。这就是欧洲封建制度的主要成因：不是"设计出来的"，而是在中央无法提供保护的前提下，地方领主用土地换取军事服役的自然反应。

\xuehai{封建契约与私人化的权力}{
  封建制度的本质是以土地（采邑）换取军事服役的个人化契约。领主将土地授予骑士（封臣），封臣向领主宣誓效忠并提供一定天数的兵役。封臣在领地上享有司法、征税和行政权力，几乎就是独立的小君主。这是一个层层下放的权力体系，而不是一个向上集中的权力体系。在这一体系中产生了一个对比清晰的差异：欧洲封建是"自下而上"的，通过私人契约层层构建；中国的分封是"自上而下"的，由天子主动授予并始终保留收回的权限。
}

\subsection{神圣罗马帝国：一个"既不神圣、也不罗马、更非帝国"的政治体}

962年，东法兰克国王奥托一世在罗马被教皇加冕为皇帝。这个后来被称为"神圣罗马帝国"的政治实体，在此后八百多年的时间里始终未能建立起和法、英两国那样的中央集权制度。1756年伏尔泰写道："这个组织既不神圣，也不罗马，更非帝国。"这句刻薄话之所以被反复引用，是因为它精准地抓住了神罗的结构性矛盾：它是一个由数百个大小邦国、自由城市和教会领地组成的松散联邦，皇帝由七个选帝侯选举产生而非世袭。

为什么神罗没能像法国那样走向中央集权？一个重要原因是帝国皇帝的选举制：每换一次皇帝，就意味着各方势力的一次重新谈判。当选者往往需要向选举者——尤其是握有重兵的大诸侯——做出大量让步以换取选票。没有世袭制的连续性，就无法积累代际优势。而在法国（以及英国），王位在同一血统内代代相传，积累的土地、制度和声望形成了不可逆的集权趋势。\textbf{选举皇帝不是民主的先声——它恰恰是中央集权失败的结构性原因。}

从查理曼到神圣罗马帝国，欧洲政治秩序最核心的特征已经显现：没有一个统一的帝国，没有一个普世的语言，没有一个可以使所有的政治碎片熔于一炉的中央权威。欧洲的分裂不是一次凡尔登条约造成的，是被继承制度、地理屏障和封建契约不断巩固、持续再生产出来的系统状态。

\subsection*{A组：内容复述与要点整理}

\begin{enumerate}[label=A\arabic*.]
  \item 查理曼为什么要在公元800年接受教皇加冕？从查理曼和教皇各自的角度分别解释。
  \item 凡尔登条约（843年）将查理曼帝国分成了哪三部分？每一部分大致对应今天的哪些国家？
  \item 法兰克人的诸子平分继承制度与中国的嫡长子继承制有什么根本差异？这一差异在一千年的时间尺度上产生了什么政治后果？
  \item 欧洲封建制度形成的主要原因是哪三方面的入侵？它和查理曼帝国的崩溃有什么因果关系？
  \item 神圣罗马帝国为什么没有像法国和英国那样走向中央集权？
\end{enumerate}

\subsection*{B组：深入分析与扩展讨论}

\begin{enumerate}[label=B\arabic*.]
  \item \textbf{比较分析题：}回顾《中国古代史》第一课中关于周代分封制的论述。比较周代分封与欧洲封建制度的异同，至少从权力来源、继承规则和中央与地方关系三个维度展开。为什么同样是以土地换忠诚，中国走向了郡县制而欧洲走向了多国体系？
  \item \textbf{历史因果分析：}843年的凡尔登条约是否"不可避免地"导致了欧洲永久性的政治分裂？反事实推理：如果查理曼只有一个继承人——欧洲的历史会有什么不同？请提出一个具体的假设并讨论其合理性。
  \item \textbf{论述题}（500-600字）：伏尔泰说神圣罗马帝国"既不神圣，也不罗马，更非帝国"。请结合神罗的选举制、多国结构和与教皇的关系，分条论证或反驳伏尔泰的论断。
\end{enumerate}

\newpage
\section{第二课\quad 皇帝与教皇：谁能任命主教？}

\subsection{一个看似技术性的问题}

今天的人们很难理解11世纪的这场冲突为什么会持续半个世纪、牵动整个西欧、甚至让一位赤脚的皇帝在雪地里站了三天三夜。问题的表面形式极其简单：地方主教应该由皇帝任命，还是由教皇任命？这个"人事任命权之争"之所以引爆，是因为教区和世俗领地在中世纪西欧是高度重叠的——一个主教不只是掌管灵魂的牧人，还管理者大片土地、拥有独立的行政权和司法权、甚至能召集自己的武装力量。在德意志，许多主教还是重要的帝国诸侯，拥有铸币权和关税征收权。

当一个君主把忠诚的亲戚或心腹安插到时富庶的教区时，他控制的是一个庞大封建领主的全部资源。教皇对此长期不满却无能为力——教皇体制本身在整个9世纪和10世纪也陷入了低谷。从公元882年到955年，二十多位教皇中有相当一部分死于非命，或被政敌囚禁，或被地方贵族操控。教皇的权威在罗马城里时常受到地方贵族家族的掣肘，更谈不上在整个基督教世界发号施令。

从10世纪开始，一场自下而上的教会改革运动从勃艮第的克吕尼修道院发起。克吕尼改革的核心诉求极其简练：教职人员的生活应当严格遵循修道院规则，反对教士结婚（当时的许多教区神父有妻室），反对买卖圣职（即"以金钱获取教会职位的腐败"）。克吕尼运动在贵族和骑士阶层中获得了大量支持，迅速扩展成为一股覆盖整个西欧的精神力量。

到了11世纪中期，克吕尼改革的精神开始影响教皇选举制度。1059年的拉特兰大公会议规定教皇由红衣主教团选举产生，不再受罗马贵族或德意志皇帝的干预。这一改革为后来格里高利七世与亨利四世的激烈对抗准备了制度的武器和意识的动员。

1075年，格里高利七世发布了《教皇敕令》，宣布教皇拥有至高无上的权威。

\subsection{从卡诺莎到沃尔姆斯：一场没有赢家的战争}

25岁的神圣罗马帝国皇帝亨利四世收到敕令后立即召集忠于他的德意志主教开会，宣布废黜格里高利七世。格里高利回应了一招中世纪政治中杀伤力最大的武器：他将亨利四世开除教籍，并宣布他的臣民的效忠誓言自动失效。

在一个大多数人真心相信被开除教籍的国王死后一定会下地狱的时代，这条命令的效果是核弹级的。德意志的贵族们——本来就不满这个年轻高傲的国王——找到了一个完美的政治借口：我们不能服从一个被开除教籍的人。大批贵族叛变，扬言要另立新君。

亨利发现如果他不低头，他的王位将在几周之内被选帝侯会议没收。1077年1月，他带着妻子和少量随从翻越了冰封的阿尔卑斯山，来到教皇当时所在的卡诺莎城堡。教皇拒绝见他。亨利在城堡门外的雪地里站了三天——赤着脚、穿着悔罪者才穿的粗毛服装。"卡诺莎之辱"成了欧洲政治史上一个标志性场景：世俗权力的最高代表跪在了宗教权力的代表面前。

\xuehai{开除教籍为什么有如此之大的政治杀伤力？}{
  在中世纪，一个人被开除教籍意味着他不能进入教堂、不能参加圣礼、不能接受临终敷油——死后将直接进入地狱。对普通信徒来说，这是终极恐怖；对国王来说，开除教籍不仅影响他个人灵魂的归宿，更是给所有想造反的封臣提供了一个完美借口。"上帝不再保佑这个国王了"——效忠誓言在对上帝的关系面前，突然变成了一纸可以随时撕毁的文书。

  格里高利七世之所以能迫使亨利低头，不是因为教皇武装比皇帝军队更强，而是因为他垄断了一项皇帝没有东西：救恩的发放权。开除教籍的威慑力恰好等于当时社会对地狱恐惧的深度。
}

但卡诺莎的戏剧性场景很容易让人误以为教皇从此凌驾于皇帝之上。事实远非如此。亨利回国后迅速集结军队击溃了他的反对者，然后立即挥师南下意大利，攻入罗马，赶走了格里高利七世。1085年，格里高利在被诺曼人从罗马救出后死于流亡之中，留下了一句广为流传的遗言："我热爱正义，憎恨不义——所以我死于流放。"

最终的和解是1122年达成的《沃尔姆斯宗教协定》。双方各退一步：教会独自拥有任命主教的权力（选举）。皇帝有权在世俗领域向新任主教授予象征世俗权力的权杖和戒指，且主教必须在皇帝的见证下方可就职。双方的妥协将宗教任命与世俗授权在程序上分割开来。谁也没有得到自己最初想要的一切——但双方都获得了可以接受的稳定局面。

\subsection{东西方的分叉：教会与国家关系的两种模式}

在拜占庭帝国，皇帝是教会的最高首脑。君士坦丁堡牧首由皇帝任命，不能独立于皇权行动。这一传统在东正教世界始终延续：教会是国家的附属机构。

在西欧，叙任权斗争的长期后果是划出了一条模糊但重要的边界：教会和世俗权威分属两个独立的权力领域，各有自己的法统源头，互为制约但互不从属。后来的思想家（从12世纪的格拉提安到16世纪的加尔文到17世纪的洛克）正是以这一"两个领域各有各的权威来源"的分裂前提为基础，发展出了更宏大的权力分散和有限政府的政治理论。

这并不是因为教皇比皇帝更明智。而是因为法兰克—德意志的皇帝从来就不够强大到完全吞并教权。拜占庭皇帝继承了完整的罗马行政机器，从一开始就能把教会纳入国家体系。西欧的帝国机器在日耳曼入侵中被打碎了，皇帝只能控制有限的领土和军队，他不得不与一个独立于他之外的组织谈判。政治制度的差异，正是在很大程度上源自双方起始力量的对比。

\subsection*{A组：内容复述与要点整理}

\begin{enumerate}[label=A\arabic*.]
  \item 11世纪教皇与皇帝之间的斗争焦点是什么？为什么教区主教的任命权如此重要？
  \item 格里高利七世用什么手段迫使亨利四世低头？开除教籍为何在中世纪具有如此强大的政治杀伤力？
  \item 1077年卡诺莎事件的结局是教皇赢了，还是皇帝赢了？简述事件后的实际走向。
  \item 《沃尔姆斯宗教协定》（1122年）是如何调解这场冲突的？
  \item 拜占庭帝国与西欧在"政教关系"上有什么区别？
\end{enumerate}

\subsection*{B组：深入分析与扩展讨论}

\begin{enumerate}[label=B\arabic*.]
  \item \textbf{制度分析：}格里高利七世和亨利四世的冲突实质上是两种合法性来源的对撞——"您被选举代表全体人民"vs"我的法典传承予我同样的权威"。你如何评价这两种主张各自的有效性和局限性？
  \item \textbf{比较视角：}拜占庭的政教合一体制维持了约一千年——西欧的政教分离却经历了频繁的内部冲突。你认为哪一种模式更有利于长期的政治稳定？哪一种模式更有利于个人自由的发展？请从社会结构和权力制衡两个角度论证。
  \item \textbf{论述题（500-600字）：}有学者认为叙任权斗争是欧洲走向近代宪政的"第一步"。你是否同意这一判断？请从这场斗争之后权力结构发生的变化展开论证：至少涉及教会权力的独立化、契约观念的形成、有限政府的早期萌芽三个要点。
\end{enumerate}

\newpage
\section{第三课\quad 西班牙：八百年收复失地与一个帝国的锻造}

\subsection{711年：一座海峡的跨越}

公元711年4月，大约七千名阿拉伯和柏柏尔士兵在北非的丹吉尔集结，在指挥官塔里克·伊本·齐亚德的率领下渡过直布罗陀海峡——后来这座海峡的名称"直布罗陀"（Jabal Tariq），在阿拉伯语中正是"塔里克之山"的意思。

西哥特王国——统治伊比利亚半岛近三百年的日耳曼政权——在当年夏天的瓜达莱特战役中遭遇了军事上的灾难。国王罗德里戈战死，王国中枢陷入权力真空。两年之内，塔里克的部队已经攻占了托莱多（西哥特首都）、科尔多瓦和塞维利亚。到718年，除了北部坎塔布连山区的一些据点外，几乎整个伊比利亚半岛都已落入穆斯林之手。

\subsection{后倭马亚王朝：一个不逊于巴格达的文明中心}

征服之后，西班牙成为倭马亚哈里发国的一个行省。750年，倭马亚王朝在远方的巴格达被阿拔斯王朝推翻，几乎全部倭马亚王室成员被屠杀——但一位年轻的王子阿卜杜勒-拉赫曼成功逃脱，辗转北非，抵达西班牙。756年，他在科尔多瓦自立为埃米尔，建立了后倭马亚王朝（756—1031年），与巴格达的阿拔斯哈里发国分庭抗礼。

鼎盛时期，科尔多瓦的人口超过四十万，是当时西欧最大的城市，也是唯一可以与君士坦丁堡和巴格达相提并论的文化中心。科尔多瓦大清真寺的规模、装饰和工程复杂程度，令造访的基督教使者目瞪口呆。后倭马亚王朝的统治者为西班牙带来了当时欧洲最高水平的农业技术——包括完善的灌溉系统、新作物品种（柑橘、水稻、茄子、菠菜、棉花和甘蔗）——以及大批翻译成阿拉伯语的希腊哲学、医学和数学著作。

\xuehai{科尔多瓦图书馆（al-Hakam II时期）}{
  后倭马亚王朝第二任哈里发哈卡姆二世（961—976年在位）投入大量资源在科尔多瓦建立了一座庞大的图书馆。据史料记载，图书馆收藏了约四十万册书——这个数字即使考虑了古代编年史的夸大，也远远超过了同时代任何一座西欧修道院的藏书。哈卡姆二世向开罗、大马士革、巴格达派出购书使者，搜寻并复制各领域的学术著作。图书馆雇佣了大批抄写员、装订匠和翻译者。

  作为对比：同时期法国巴黎的西岱岛上，所谓的"宫廷学校"的藏书大约只有数百册之谱。百年翻译运动的主要工作，在很大程度上也依赖于分散在西班牙诸图书馆中的阿拉伯文稿库。
}

\subsection{基督教王国的推进}

在科尔多瓦文明辐射着地中海西部的几个世纪里，伊比利亚半岛北部山区的几个基督教小王国——阿斯图里亚斯、莱昂、纳瓦拉、卡斯蒂利亚、阿拉贡——正在缓慢而坚定地扩张。这个过程被称为"收复失地运动"（Reconquista），但它的面貌并非常年战争：在几个世纪的时间里，前后有数十年的和平期，大量的交流和贸易在王国与埃米尔国之间进行。

\textbf{988年，科尔多瓦的内部崩溃开启了形势转换。} 后倭马亚王朝末代哈里发于1031年被废黜，哈里发国解体为二十多个互相争斗的小王国。这给半岛北部基督教王国提供了利用分裂局势、逐个削弱的窗口。此前几个世纪的不定期进攻，在这一时期转换为系统性的向南推进。

1085年，卡斯蒂利亚国王阿方索六世攻克托莱多——这是基督教王国从穆斯林手中夺取的最大城市。托莱多的失陷震动了穆斯林世界，北非的穆拉比特王朝应西班牙穆斯林统治者的请求渡海援战。此后两个多世纪中，新的北非王朝（穆瓦希德王朝）与基督教王国在半岛上反复拉锯。

关键时刻出现在1212年的纳瓦斯·德·托洛萨战役。卡斯蒂利亚国王阿方索八世在教皇号召下联合了阿拉贡、纳瓦拉和葡萄牙的军队——外加一群主动参战的法国十字军骑士——在安达卢西亚腹地与穆瓦希德哈里发的主力大军正面对决。联军获胜，穆瓦希德王朝在伊比利亚半岛的军事实力遭到毁灭性打击。此后半个世纪中，安达卢西亚的大部分地区迅速落入基督教王国之手。

到13世纪末，穆斯林在伊比利亚半岛的控制区已经缩小到只剩格拉纳达王国（今天安达卢西亚最南端的一小块多山地区）。格拉纳达的纳斯里王朝得以在此后的两百年间继续存续，主要原因在于它灵活的外交手腕：它向卡斯蒂利亚称臣纳贡，同时在北非的梅里尼德王朝和半岛上的几个基督教王国之间保持微妙的平衡。格拉纳达的阿尔罕布拉宫——纳斯里王朝宫殿群——在建筑装饰上达到了伊斯兰西班牙艺术的最高水准。它的墙壁上反复出现的一句铭文——"万物非主，唯有安拉"——被雕刻在石膏和木雕的精致文饰中，构成了一种文字与空间合一的美学表达。

\xuehai{收复失地运动为何持续了接近八百年而不是一次性的军事征服？}{
  从718年科瓦东加战役算起，到1492年格拉纳达陷落，前后持续了七百七十四年。这一时期如此漫长，关键约束在于中世纪的技术条件：崎岖的山地地形限制了后勤补给线路，无法维持征服者长驱直入；北部基督教诸国的人口基数长期不足以支撑决定性歼灭战；南方的穆斯林政权在受到严重威胁时可以召唤北非的援军。八百年间，和平期的长度加起来远超过战争期。
}

\subsection{1492：西班牙从边缘走向舞台中央}

1469年，阿拉贡王储费尔南多与卡斯蒂利亚女王伊莎贝拉联姻，两个最大的基督教王国实现了合治。此后他们集中力量攻打半岛上最后一块穆斯林飞地——格拉纳达王国。经过十年断断续续的围城和山地作战，1492年1月2日，格拉纳达的末代苏丹将城门钥匙交给了费尔南多和伊莎贝拉。收复失地运动终结。

同一年三月，费尔南多和伊莎贝拉签署了《阿尔罕布拉法令》：境内所有犹太人在四个月内要么改信天主教，要么离开西班牙。据估算，约十万到二十万犹太人被驱逐出境，其中包含大量医生、天文学者、商人和银行家。1502年西班牙对境内残存的穆斯林（穆德哈尔人）下达了类似的强制改信或驱逐令。到17世纪初，西班牙又对"莫里斯科人"实施了大规模驱逐，总数约三十万人。这些驱逐使西班牙在宗教上获得了统一，但代价是长期的经济和人才流失。

同一年八月，一个热那亚航海家接受了费尔南多和伊莎贝拉的资助，率领三艘小船从帕洛斯港出发向西横渡大西洋。他的名字叫克里斯托弗·哥伦布。

1492年的三件事情——终结收复失地、驱逐犹太人、资助哥伦布——不是时间上的巧合。八百年边疆战争塑造的西班牙：以天主教为精神核心，以国王为政治中心，以武力和征服为国家经验。费尔南多和伊莎贝拉在统一半岛后面临一个严峻的问题：他们解散了大量的军事组织，手中握着一支身经百战但无所事事的骑士队伍。把这些能量导向大西洋几乎是顺理成章的——在伊比利亚半岛上没有人可以再打了。从地球的这一端到另一端，这笔账是同一个出纳员结算的。

\subsection*{A组：内容复述与要点整理}

\begin{enumerate}[label=A\arabic*.]
  \item 711年穆斯林征服伊比利亚半岛的经过是怎样的？
  \item 科尔多瓦后倭马亚王朝在文化上的主要成就有哪些？
  \item "收复失地运动"持续了约多少年？在前几个世纪中进展缓慢的原因是什么？
  \item 1492年同时发生的三件大事是什么？它们之间有没有内在联系？
  \item 费尔南多和伊莎贝拉为什么要在统一半岛之后支持哥伦布的远航？
\end{enumerate}

\subsection*{B组：深入分析与扩展讨论}

\begin{enumerate}[label=B\arabic*.]
  \item \textbf{比较分析：}科尔多瓦哈里发国同时期的"阿拔斯百年翻译运动"与"伊比利亚翻译浪潮"有什么异同？为什么古希腊和罗马的学术遗产在欧洲边缘的西班牙反而保存得比罗马城本身更为完整？
  \item \textbf{因果链分析：}西班牙收复失地运动塑造了什么样的社会结构？请列出至少三条它对外投射力量的"遗产"（军事制度、宗教政策、国家组织方式），并联系查理五世和菲利普二世的帝国战略加以论证。
  \item \textbf{论述题}（500-600字）：1492年西班牙驱逐犹太人是否是必须进行的历史清算？请从国家统一、财政收入、智力资本外流和国际声誉四个角度评估该政策的得失。
\end{enumerate}

\newpage
\section{第四课\quad 英格兰：法律如何一步一步管住了国王}

\subsection{诺曼底的海上跳跃}

1066年1月，英格兰国王忏悔者爱德华去世，无子嗣。三位声称有继承权的竞争者中，诺曼底公爵威廉最有准备。他在当年秋天率领约七千人的部队和一支庞大的补给舰队渡过英吉利海峡，在佩文西登陆。10月14日在黑斯廷斯战役中击败并杀死了英格兰国王哈罗德。圣诞节当天，威廉在伦敦威斯敏斯特大教堂加冕为英格兰国王。

征服者威廉做的第一件事是一场大规模的行政抄写——他派遣专员走遍英格兰每一块土地，登记每一座庄园、每一片牧场、每一条河流的渡口，甚至每一头猪。调查的结果被编成两卷《末日审判书》（Domesday Book），意为"最后的审判之书"，因为它对财产和权利的记录被认为像最后的审判一样无可回避。这一调查的目的是税收精确化：国王现在知道全国有多少可征税的财产和它们的所有者。与此相比，同一时期的法国国王连巴黎附近有多少伯爵领地在技术上属于王室的直接管辖都说不太清。

\xuehai{中央集权的"天然空间"：英格兰和法兰西的差异是如何产生的？}{
  诺曼征服后的英格兰之所以迅速建立了全国性的土地调查和统一的行政体系，关键原因是这块土地的面积和形状适合行政管控。英格兰的面积约13万平方公里，在11世纪骑马信使的极限传递半径内。巴黎盆地周边的法兰西岛面积不足英格兰的十分之一。英格兰的统一行政能力并非"天赋"，而是被国土面积和地形的适中程度所赋予的。
}

威廉把全国封建主从土地到军役全部编织进了一张以国王为顶点的网络中。英国的封建制度和法国的封建制度从一开始就有差异：在法国，"我的附庸的附庸不是我的附庸"——国王无法直接指挥大贵族手下的骑士；在英格兰，威廉强行要求全体封臣——不论大小——直接向他宣誓效忠。

征服者威廉建立的高效中央机器为金雀花王朝的进一步扩张提供了基础平台。1154年登基的亨利二世不但是英格兰国王，还通过联姻继承了诺曼底、安茹、阿基坦——他实际控制的法国西部领土面积超过了法国国王自身的王室领地。这一时期被称为"安茹帝国"。亨利二世还系统性地改革了英格兰的司法体系。他将国王法庭的巡回审判制度固定在习惯法中——在地方治安法庭和国王的中央法庭之间建立起了一套后续数百年的连接制度。陪审团被广泛运用于解决私人和产权纠纷。这一时期形成的英格兰普通法传统，与欧洲大陆系统化的罗马法传统之间开启了一条平行的、直到今天仍然保持独立的发展路径。

安茹帝国在亨利二世的儿子约翰手中迅速瓦解。1204年，约翰对法战败，诺曼底被法王腓力·奥古斯特吞并。王室收入的一半以上的领土税基全部丧失。约翰随后与教皇英诺森三世的冲突导致英格兰被置于禁止举行弥撒的禁令下长达数年。

\subsection{1215年兰尼米德：大宪章与它的语境}

诺曼和金雀花王朝在随后一个半世纪中建立了英格兰历史上最高效的封建行政机器。但这部机器在林子里最凶险的一次考验来自于一个非常糟糕的国王。约翰王在1199年登基后做的事情几乎每一件都对国家造成了伤害：1204年对法战争大败，丢掉了诺曼底（英格兰王室在法国北部最重要的领土），他因此获得了"失地王"的绰号；他与教皇英诺森三世的冲突导致英格兰被置于禁令之下长达六年，全国教堂关闭，不得举行弥撒。

1215年6月15日，一群全副武装的贵族在温莎附近的兰尼米德草地上截住了约翰王，逼迫他在一份长达六十三条的法律文件上盖了国王的玉玺。这份文件是《大宪章》。

必须强调的是：《大宪章》不是一部关于"人类自由"的普遍宣言。它的核心诉求是保护贵族和教会的既得利益不受国王的任意侵犯。其中最关键的条款是第三十九条——"任何自由人，未经同等地位者依法审判或国家法律裁决，不得被逮捕、监禁、剥夺财产、流放或以任何方式毁灭。" 虽然"自由人"在1215年的英格兰仅占总人口的一小部分（大量的农奴不在此列），但这一条款第一次以成文法的形式写下了"国王不能仅凭个人好恶就抓人"的原则。

\xuehai{大宪章为什么在后世被反复援引？}{
  大宪章在整个中世纪的实际效力波动很大。它在1215年签署后很快就被教皇宣布无效。但它被后续的君主在不同时期为了换取征税权而重新确认了约三十次。每一次重新确认都更新了它在法律实践中的参考地位。17世纪的英国律师和议会反对派在反对查理一世的专制统治时，正是从大宪章中找到了"国王必须遵守法律"的历史先例。一部在中世纪主要保护贵族特权的法律，变成了为议会主权原则提供正当性的历史基础。
}

\subsection{从内战到复辟再到光荣革命：三步不是一步}

从大宪章到议会主权的确立，经历了四个多世纪的反复。13世纪后期，爱德华一世为了筹集对威尔士和苏格兰的战争经费，频繁召集贵族、教士和平民代表开会讨论税收——这个为了"搞钱"而召集的会议逐渐演变成议会。此后数百年的核心情节是：国王需要钱→议会要求权利→国王做出让步→议会批准税收。

查理一世在1628年接受了《权利请愿书》，承诺未经议会同意不得征税，但他随后在1629年解散议会并尝试了十一年无议会统治。1642年，议会与国王之间的冲突从议事厅转移到战场，内战爆发。1649年，奥利弗·克伦威尔领导的议会军获胜，查理一世被作为"暴君、叛徒和国家的敌人"受审并处决。英格兰成为共和国。

共和国在克伦威尔死后迅速陷入军事统治的无序状态。1660年，议会主动邀请查理一世的儿子查理二世从流亡地法国回来恢复君主制——反对专制不等于反对君主制本身，英国人在没有国王的军事独裁中学到了这个教训。

1685年，查理二世去世，他的弟弟詹姆斯二世即位。詹姆斯公开信奉天主教，在议会中引起了道义层面的普遍恐惧。1688年，议会中的两党领袖秘密邀请詹姆斯二世的女儿玛丽和她的丈夫——荷兰的奥兰治亲王威廉——带兵介入。威廉的军队登陆英格兰时几乎没有遇到抵抗，詹姆斯逃往法国。议会宣布他自动退位，由威廉和玛丽共同加冕为新国王。

1689年的《权利法案》是这次权力交接的最终法律文件：国王未经议会同意不得征税、不得在和平时期维持常备军、议员在议会中享有言论自由。但所有这些条款都只是框架，而非终点——整个18世纪，国王仍然通过内阁任命、选举操控和议会贿赂实质性地影响着立法和行政。议会真正成为政治生活的中心、首相取代国王成为行政首脑，那是19世纪中期经过多次选举改革后的结果。英格兰的宪法不是在1688年就完工的。它是在每一个事件为下一个事件的起点进行修补的过程中累积形成的。

\xuehai{波兰的"自由否决权"——制约权力走到极限就是亡国}{
  当英国正在逐步探索约束王权的制度手段的同时，波兰-立陶宛联邦在另一个方向上进行了截然不同的实验。从17世纪中期开始，波兰的贵族议会（瑟姆）可以凭借任何一位议员的"自由否决权"废除整个议会的全部决议——完全不需要理由。这一制度设计的初衷是保护少数派贵族不被多数派压迫——初衷和英国的"国王不可为所欲为"有相似之处。

  实际结果是灾难性的。1652年至1764年间，波兰议会召开了71次，超过一半（42次）被单个议员的否决票直接瘫痪。任何一项行政财政改革或军事动员立法都会在议会中陷入停滞。到18世纪晚期，波兰被其三个有绝对王权的邻居——俄国、普鲁士和奥地利——三次瓜分。1795年，波兰从欧洲地图上消失了。有效的权力制约需要确保制度不被少数人的一票否决权彻底锁死。英国的议会最终选择了多数决规则，波兰选择了全体一致。
}

\subsection*{A组：内容复述与要点整理}

\begin{enumerate}[label=A\arabic*.]
  \item 征服者威廉在英格兰建立了什么样的统治体系？《末日审判书》调查的最直接目的什么？
  \item 1215年《大宪章》诞生的直接原因是什么？它的第三九条写下了什么原则？
  \item 英国议会是怎么从"国王的筹款会议"演变而来的？
  \item 1642-1649年英国内战的结果如何？为什么1660年议会又请回了国王？
  \item 1688年光荣革命后通过的《权利法案》主要内容是什么？
\end{enumerate}

\subsection*{B组：深入分析与扩展讨论}

\begin{enumerate}[label=B\arabic*.]
  \item \textbf{比较分析：}波兰的"自由否决权"和英国的议会多数决——同样是约束国王权力的工具，为什么一个最终导致国家亡国而另一个却奠定了制度化的政治协商？是哪一条关键的规则设计差异造成了这种分化？
  \item \textbf{因果链分析：}英国的"君主立宪"制度是在哪些关键历史节点的推动下逐步成型的？请在大宪章、内战、复辟、光荣革命四个节点之间，补充必要的过渡（如爱德华一世的模范议会）、倒退（如十一年暴政）、反复（如共和国失败）环节，绘制一幅包含正面反馈和负面反馈节点的因果图。
  \item \textbf{论述题}（500-600字）：光荣革命后执行的政治权力格局是否符合《权利法案》字面上的承诺？请结合18世纪英国国王通过任命首相、操控选举和财政收买干预议会的事实，评价"1688年确立的是议会至上"这一论断是否经得起检验。
\end{enumerate}

\newpage
\section{第五课\quad 百年战争：英法恩怨与民族国家的诞生}

\subsection{一桩婚姻衍生了四百年的战争}

1066年诺曼征服的后遗症在英法关系中制造了一个结构性的矛盾：英格兰的国王同时也是法国的封臣——作为诺曼底公爵，他向法国国王行封臣礼。当你同时是一个主权国家的元首和另一个主权国家的封臣时，你的外交和军事行动迟早会陷入无法调和的冲突。

通过联姻、继承和战争，金雀花王朝在12世纪中控制了诺曼底之外的大片法国西部和西南部领地——安茹、曼恩、普瓦图、阿基坦。到1154年亨利二世即位时，他统辖的法国领土面积比法国国王本人直接控制的王室领地大好几倍。法国国王不能容忍这一局面，但几代人以内的武力手段都不足以将其全面驱逐。两个国家之间的领土争端被一个王朝法理上的双重身份积压进了骨子里。

\subsection{一场打了一百多年的战役集合}

1337年，法国国王腓力六世宣布没收英王在法国的领地，英国国王爱德华三世则宣布自己对法国王位拥有继承权——他是已故法王腓力四世的外孙。战争爆发。

百年战争的第一个阶段（1337-1360）由英国取得了一边倒的军事胜利。关键战役是1346年的克雷西战役和1356年的普瓦捷战役。在这两场战役中，英国的长弓手以少胜多地击溃了法国重装骑兵的多次冲锋。英格兰长弓手是训练有素的职业步兵，能在每分钟射出十到十二支穿甲箭，有效射程超过两百米——远大于弩的射速且优于弩的射程。法国贵族骑士穿着重甲骑着高头大马冲锋时，面对的是从二百米外持续不断、密集覆盖射来的穿甲箭雨。马匹被射翻、骑士倒地、笨重的铠甲使他们无法自行站立。在克雷西和普瓦捷，法国最精锐的贵族阶层的战斗部队几乎被成批杀伤。

这迫使法国在1360年签署了极其屈辱的《布勒丁尼和约》，承认了英国对法国西南部和加莱海峡的巨大领土所有权。和约维持了不到十年，法国的新一轮攻势持续到1380年代，大部分失地得到了恢复。此后进入二十年的休战期。

第二个致命打击发生在1415年。英国国王亨利五世在阿金库尔战役中再次以悬殊的兵力对比击溃了法国军队——长弓再一次扮演了决定性角色。亨利五世随后占领了诺曼底，并与勃艮第公爵结盟——勃艮第是当时法国最为独立的强势封臣，在英法战争中公开站在英国一侧。1419年亨利五世兵临巴黎城下，法国王太子被迫退到卢瓦尔河以南，英法勃艮第三方签署了《特鲁瓦条约》，条约规定亨利五世及其后代将在法国国王查理六世去世后继承法国王位——如果条约完全执行，英法两国将在一位共同的君主统治下合并为一个跨海峡的联合王国。这个前景在英国政局变动和战场上持续的变化面前最终没能实现。

\subsection{被战争碾碎的底层}

百年战争的负担不仅落在贵族和骑士身上，农村中的普通农民承担了主要的赋税、征粮和兵役。战争最残酷的一种战术被称为"敕令远征"或"劫掠行军"——英军指挥官通过有组织地烧毁村庄和农田来剥夺法国领主的税基、削弱对手的经济基础。在克雷西战役之后，爱德华三世的儿子"黑太子"爱德华对法国南部发动了著名的"劫掠行军"，从波尔多出发直抵地中海沿岸再折返，一路上农田和村落被全面破坏。这种"通过消耗农民的财产来削弱贵族税基"的战略对农村居民的打击要比对城堡中领主的打击大得多。

1358年爆发的"札克雷"（来自农民绰号"雅克好人的造反"）是法国北部农民对被战争掏空的生存条件的大规模暴力反抗。起义在几个星期内席卷了巴黎周边和皮卡第地区，农民武装攻占了多座城堡并杀害了数名当地领主。法国的贵族阶层在震惊之余迅速集结兵力，在几场战役中镇压了起义，随后进行了惨烈的报复。札克雷被镇压后，法国北部的社会经济需要数十年才能恢复到战前水平。

这种贵族之间的王朝战争由底层社会中最脆弱的群体支付最直接的代价的模式，在整个欧洲的中世纪晚期至近代早期的战争史中是持续的结构性特征。

\subsection{贞德：士气转化为了胜利}

到1429年，法军已经溃败到了岌岌可危的底线上。巴黎和北部大片地区被英军和勃艮第军队占领；法国王太子查理退到卢瓦尔河以南，连登基加冕的仪式都无法举行（加冕地兰斯大教堂在英军控制的区域）。奥尔良是卢瓦尔河上的防线缺口——如果奥尔良陷落，整个法国南部将门户大开，王太子逃亡或投降的选项会摆在桌上。

就在这个几乎所有人都失去了希望的时间点上，一个叫贞德的17岁农家女走进了王太子在希农的城堡，声称她听到了来自上帝的召唤：她要解救奥尔良，护送查理前往兰斯加冕。贞德不识一个字，没有接受过任何军事训练——但她在战场上表现了一种当时法国军队已经丧失了的东西：毫不动摇的进攻意志。

她并不是一位冲锋陷阵的骑士。她几乎从不直接参与战斗。她的作用是站在军旗前，让每一个人都看到她——然后她用剑背驱赶溃退的士兵，逼着他们回头继续进攻。在此之前，奥尔良城内的法军已经被围困了七个多月，龟缩在城墙后面不敢出战。贞德到达后九天，奥尔良围困解除。

\xuehai{长弓、火药与封建骑士的终结}{
  阿金库尔战役中英军长弓手的决定性表现，揭示了封建军事制度的一个根本缺陷：以贵族骑士为核心的重骑兵在面对有组织的职业步兵和远程投射火力时极其脆弱。克雷西、普瓦捷和阿金库尔三次大屠杀实质上消灭了法国贵族骑士阶层中相当大比例的世代军人。

  然而为封建军事制度敲响丧钟的还不是长弓。到15世纪晚期，火药武器————火炮和火绳枪————开始在战场上扮演角色。法国在1453年（百年战争结束的那一年）已经拥有欧洲最强大的火炮装备。火炮可以击穿城堡的石墙——这意味着领主的防御堡垒第一次不再绝对安全。国王的炮队可以轰开任何拒绝缴纳新税的贵族堡垒。这一技术变革深刻地改变了权力分配：国王对地方领主的军事优势在十五世纪内变成了不可逆转的定局。
}

1430年5月，贞德在一次小规模战斗中被勃艮第人俘虏。他们随即把她卖给了英国人。英国人组织了由巴黎大学神学教授组成的宗教法庭，起诉贞德为异端分子和女巫。经过数月的审讯，她被判处火刑。1431年5月30日，在鲁昂老市场广场，贞德被公开烧死。她当时未满十九岁。

1453年，战争结束。英格兰失去了除加莱港之外在法兰西土地上的全部领地。法兰西完成了统一，王权空前强大，近代法国民族国家的轮廓已经划出。

\xuehai{查理七世的军事遗产：法兰西常备军制度的确立}{
  百年战争的一项重大制度性后果是法国建立并维持了欧洲第一支常备军。1439年的《奥尔良法令》禁止了领主私人之间的战争，赋予国王对军事力量的绝对垄断权。国王招募的"敕令连队"由国王直接支付军饷，不再依赖贵族封臣的封建兵役。这是中世纪末期到近代早期西欧军事制度最关键的转折——从私人契约武装转向国家垂直控制武装。

  英格兰走向常备军的时间要晚得多。查理二世在复辟后尝试建立一支小规模的常备军，但议会出于对国王利用军队推行专制的恐惧而严格限制其规模。这一制度的差异使得法国在国际关系中获得了军事动员效率上的不对称优势。
}

\subsection{一场战争怎么造出了两个民族}

百年战争最重要的遗产不是领土的得失，而是认同的边界。在战争之前，一个普通的法国农民可能从未想过自己是否属于一个叫"法兰西"的国家。他的效忠对象是领地的领主：至于领主效忠于英格兰国王还是法兰西国王——那是领主自己的事。但当一个农民在一百年的时间里反复经历英军的焚烧、掳掠和屠杀，他的自我意识开始发生转变。对面的那些人不再是"另一个领主的人"，他们是"英国人"。"我们"和"他们"的分界线，在一百二十年的战火中被清晰地从潜意识里拉了出来。百年战争不是民族主义的原因——但它是一块把语言、领地和历史暴力粘连在一起的焊接剂。

\subsection*{A组：内容复述与要点整理}

\begin{enumerate}[label=A\arabic*.]
  \item 英法百年战争的根本起因是什么？金雀花王朝的领土问题是如何形成的？
  \item 克雷西和普瓦捷战役中，英军以少胜多击败法军的关键技术因素是什么？
  \item 贞德在奥尔良解围中发挥了什么样的作用？
  \item 贞德被俘后，为什么英方组织了由神学家组成的宗教法庭来审判她？
  \item 百年战争对于法国和英国各自的民族认同产生了什么影响？
\end{enumerate}

\subsection*{B组：深入分析与扩展讨论}

\begin{enumerate}[label=B\arabic*.]
  \item \textbf{制度分析：}英格兰长弓在百年战争中暴露了以贵族重骑兵为核心的封建军事制度的哪些结构性缺陷？请从兵源储备、作战训练时间和装备成本三个维度分析为何长弓能够碾压重骑兵。带出火药武器对城堡防御的技术性淘汰这一后续影响。
  \item \textbf{比较人物分析：}贞德和《中国古代史》中的花木兰都是"女性参军"的传说/历史人物。请比较两人在各自国家集体记忆中所占据的地位有何不同——影响各自形象塑造的因素（真实史料留存度、民族危机形成的时代、当代政治宣传需求等）有哪些？
  \item \textbf{论述题}（500-600字）：百年战争是否可以被视为"欧洲民族国家的起源"？请从战争对领土边界（英格兰放弃大陆领土）、战争对军事制度（骑士地位的衰落与常备军的兴起）和政治整合（法国王室权力的集中）三个层面展开论证。
\end{enumerate}

\newpage
\section{第六课\quad 三十年战争：主权国家的规则是怎样定下来的}

\subsection{宗教改革留下的火药桶}

1517年马丁·路德在维滕堡教堂大门上钉出《九十五条论纲》之后，西欧的宗教格局从大致统一的天主教世界迅速分裂。1555年的《奥格斯堡和约》确立了"教随国定"（Cuius regio, eius religio）的原则：各邦的统治者可以决定其领地内臣民的宗教信仰。这一安排暂时平息了德意志地区的宗教冲突，但它没有解决任何深层矛盾。

和约只承认路德宗和天主教为合法信仰——加尔文宗未被纳入。与此同时，德意志南部的几个大主教辖区和修道院领地的控制权在天主教与新教诸侯之间持续拉锯。当哈布斯堡家族——同时掌握着神圣罗马帝国皇位、西班牙王位、尼德兰和意大利部分领土——试图在天主教占据优势的地区扩大势力、在边疆地区推动或强制反宗教改革时，等待的已经不是一个可以谈判的局面。

\xuehai{哈布斯堡包围圈：法国对外政策的核心焦虑}{
  到16世纪晚期，哈布斯堡家族控制下的领土包括了西班牙、葡萄牙（1580年起联合）、荷兰（当时还在西班牙统治下）、弗朗什-孔泰、米兰、那不勒斯、西西里和神圣罗马帝国的皇位。这从南、北、东三个方向围绕着法国形成了一道包围圈。法国决策者无论从哪一个方向往边境以外看，看见的都是哈布斯堡的堡垒和领地。因此，从16世纪到18世纪中期的法国对外政策有一条贯穿始终的核心逻辑：在一切可能的场合削弱哈布斯堡的势力。
}

\subsection{1618：布拉格的那座窗户}

1618年5月23日，一群波希米亚新教贵族冲进布拉格城堡，将两位皇帝派驻的总督和他们的秘书从城堡窗口扔进了护城河。三个人掉进了窗下的垃圾堆——捡回了命。布拉格抛窗事件在波希米亚引发了贵族起义，并且迅速将欧洲大陆卷入了后来被称为"三十年战争"（1618—1648）的全面性冲突。

战争上半场（1618-1629）相当简明地描述为"哈布斯堡占了上风"：波希米亚起义在1620年的白山战役中被迅速镇压；丹麦国王克里斯蒂安四世介入后也在1626年被皇帝雇佣的瓦伦斯坦将军击溃；到1629年，皇帝斐迪南二世发布了《归还敕令》，要求所有自1552年以来被新教诸侯占用的教会土地归还给天主教会。哈布斯堡似乎即将在德意志恢复天主教的主导地位。

但1630年改变了战争的走向。瑞典国王古斯塔夫二世·阿道夫在德意志北部登陆，在一系列战役中展示了经过他改革的瑞典军队的惊人战斗力——瑞典是欧洲最早全面装备火绳枪和火炮并将其与步兵机动战术相结合的国家。1631年在布赖滕费尔德击败了皇帝的天主教联军，1632年在吕岑战役中击溃了瓦伦斯坦的军队——尽管古斯塔夫本人在吕岑战役中阵亡。

1635年，天主教法国对新教瑞典的援助从秘密资助变为直接出兵参战。法国本土的安全虽然没有受到波希米亚或巴伐利亚的威胁——但黎塞留的判断基础是：哈布斯堡的势力一旦在德意志巩固，法国将面临被全面封锁的局面。法国的加入使得战争从一场德意志的宗教内战彻底转变为全欧范围的权力平衡战。

1643年的罗克鲁瓦战役标志着西班牙陆军威风的终结。法国年轻的将领孔代亲王在罗克鲁瓦（今法国北部）以法军重骑兵和步兵的联合冲击击溃了南尼德兰的西班牙军队——在此之前的十几场战役中，西班牙的"大方阵"步兵战术在欧洲战场上未曾遭受过同等规模的失败。罗克鲁瓦之后，战场优势明显地向法国一侧倾斜。与此同时，瑞典军队在波罗的海南岸持续发动攻势，占领了大片北德意志领土。哈布斯堡皇帝斐迪南三世在1640年代中期已经无力将所有战线统一指挥，和谈的军事基础初步形成。

\xuehai{黎塞留的"国家理由"（raison d'état）及其影响}{
  红衣主教黎塞留是路易十三的首相，也是天主教会的枢机主教——但他做出的战略决策是在财政上支持新教的瑞典军队去对抗日耳曼的天主教皇帝。对黎塞留来说，\textbf{国家利益高于宗教信仰}，这是一个决定性的思想——与现代国际关系中"主权至上"的原则同属一条谱系。

  黎塞留的《政治遗嘱》中写道：\textbf{"对人来说，保存自身是首要法则；对王国来说，保存国家是首要法则。"} 这标志着从中世纪的"统一基督教共和国"观念向"民族国家利益"的正式转换。三十年战争的后半段，宗教标签几乎完全退化为外交形式，实质上是哈布斯堡家族与波旁家族之间的争霸战。
}

\subsection{威斯特伐利亚：旧秩序的死亡与新规则的确立}

1648年签订的《威斯特伐利亚和约》终结了三十年战争。它的深远影响不在于单一宗教问题的解决——事实上它在德意志确认了天主教、路德宗和加尔文宗三种信仰的合法并存——而在于它建立了一套全新的国际关系规则：

第一，主权平等。每个缔约国的领土和管辖权不受外来干涉。在威斯特伐利亚之前，欧洲的外交逻辑中存在着一种"更高的正当性"——教皇或皇帝有权干预一个"走上错误道路"的王国。威斯特伐利亚的原则宣告了这种正当性的终结：外部势力不能以宗教或某种更为崇高的理由干预另一个主权国家的内部事务。

第二，帝国名存实亡。和约承认神圣罗马帝国境内的各邦享有独立的外交权和结盟权。帝国变成了一个空壳，此后一百五十年里持续存在，但渐渐转型为地区政治团体的协商框架，而非一个真正行使治理权力的机构。

第三，多国体系取代普世帝国。欧洲从此不再存在一个试图协调整片大陆的权威（不论是教皇还是皇帝）。取而代之的是若干彼此承认独立地位、通过外交和条约相互打交道的"主权国家"。

\xuehai{威斯特伐利亚原则的当代回响}{
  1945年《联合国宪章》第二条第7款规定的"不干涉本质上属于任何国家国内管辖之事件"的原则，其文本来源可以追溯到威斯特伐利亚和约。2008年俄格战争后俄罗斯承认南奥塞梯独立时援引的正是科索沃先例——实际上是在威斯特伐利亚主权话语体系中的争论。2022年俄乌战争爆发后，双方各自援引威斯特伐利亚体系的不同原则：俄罗斯援引"不得不干涉邻国的内部事务"以回应北约的持续东扩；乌克兰援引"主权国家边界不容外部武力改变"这一自1648年以来形成的国际法基线。
}

\subsection*{A组：内容复述与要点整理}

\begin{enumerate}[label=A\arabic*.]
  \item 三十年战争爆发的直接导火索是什么？
  \item 丹麦和瑞典先后介入德意志战争，各自的动机有什么不同？
  \item 为什么天主教法国在1635年站在新教瑞典一边对天主教皇帝宣战？
  \item 《威斯特伐利亚和约》对国际关系的最重要贡献是什么？
  \item 威斯特伐利亚规则对神圣罗马帝国的影响是什么？
\end{enumerate}

\subsection*{B组：深入分析与扩展讨论}

\begin{enumerate}[label=B\arabic*.]
  \item \textbf{制度变迁分析：}三十年战争期间"宗教动机"是如何一步步退居次位、让位于"国家理由"的？请以黎塞留、古斯塔夫二世·阿道夫和瓦伦斯坦三个关键人物为例说明各自决策过程中宗教因素和权力因素的相对权重变化。
  \item \textbf{概念分析题：}什么是"主权"？威斯特伐利亚体系确立的"主权"概念在1648年首次做出的核心界定在当代受到了哪些挑战？请至少列举气候变化、人权干涉和跨国资本流动三个领域加以说明。
  \item \textbf{论述题}（500-600字）：威斯特伐利亚和约是一次成功的和平，还是一份在灾难中被迫接受的权力重组协议？请以德意志的人口损失、领土再分配结果和国际体系的未来走向为论据展开论述。
\end{enumerate}

\newpage
\section{第七课\quad 理性革命：从宇宙图景到政治设计}

\subsection{哥白尼与牛顿：宇宙的重新定位}

1543年，一位波兰的教会法学者在临终前出版了一部挑战了世界观的著作。尼古拉·哥白尼在《天体运行论》中提出：太阳并非绕着地球转——恰恰相反，地球和其他行星都在围绕太阳运行。在接下来的近一个世纪里，这一论点获得了越来越多天文学观测证据的支持。1609年，帕多瓦大学数学家伽利略用自制的望远镜观测到了月球表面的环形山、木星的四颗卫星和金星的相位变化。这些观测结果与托勒密的地心说不可调和，却与哥白尼的日心说完全吻合。

伽利略将他的发现整理后出版了《关于托勒密和哥白尼两大世界体系的对话》。他在书中借一个代表教皇观点的人物"辛普利乔"（Simplicio，意为"单纯的人"）之口陈述了反对日心说的论点。1633年，宗教裁判所判定伽利略"有强烈异端嫌疑"，判令他公开放弃日心说并判处软禁。据传，伽利略在宣判后低声说了一句"但它确实在动"——这句话的可靠性存在争议，但它高度概括了他与教会之间冲突的根本性质：他指控的结论是基于数学计算和系统观测整理出来的数据。教会指控他违反了《圣经》的字面含义。

1687年，伊萨克·牛顿出版了《自然哲学的数学原理》。在这部著作中，牛顿用万有引力定律和三大运动定律统一解释了天体运动和地面物体运动的因果机制。月球绕地球运行的轨道运动规则在地球上一颗苹果落地的斜线运动都归因于同一套力场方程。牛顿并不是一个"用科学取代宗教"的启蒙斗士——他花费了至少二十年的时间研究《圣经》的年代学和启示录预言——但他的物理学产生了一个连他自己也未曾预见的思想冲击：如果宇宙是一台按照精密数学规则自主运行的机器，那么上帝在这台机器中担任的是什么角色？一个需要持续介入这个世界来维持秩序和审判、还是只负责上好了发条就开始运转的无为管理者？

\xuehai{机械宇宙论的兴起}{
  牛顿物理学将大自然的运作描述为一套可以通过数学方法加以描述和预测的因果体系。这一观念对启蒙思想家的吸引力远远超出了天文学的边界：如果大自然遵循可以被发现的规律——那么人类社会的运作是否也应该存在着类似的规律？欧洲的统治者是否也能够被一套合理的制度加以约束而被纳入某种预先确定的总体秩序中？

  这一对"社会物理学"的信念——由约翰·洛克、伏尔泰和孟德斯鸠分别以不同的方式在不同领域展开——构成了启蒙运动社会政治思想的深层理论基础。这些政治思想家用牛顿的严谨推理来考察法律、宗教权威和政治权力的结构：不再以习俗为依据，而是以理性标准进行判断。

\subsection{知识的方法论革命：培根、笛卡尔与新的求知路径}

在启蒙运动的政治方案开始涌现之前，一场更为基础的方法论革命已经在发生：人们改变了对"如何获得可靠知识"的认知。

弗兰西斯·培根（1561—1626）在其《新工具》中批评了经院哲学对亚里士多德著作的权威性依赖，提出知识必须从系统收集事实开始，通过归纳法而非先验推理获取结论。培根的名言"知识就是力量"不仅是一句关于认知的格言，更代表着对学术研究的社会功能的重新界定——知识不应该被困在修道院和大学课堂上，它应该实实在在地推动技术进步和物质条件的改善。

勒内·笛卡尔（1596—1650）走了另一条路。在他的《方法论》中，笛卡尔提出从普遍怀疑出发——怀疑一切感官提供的信息、怀疑数学是否确实可靠、甚至怀疑外部世界是否真实存在。在这种激进的怀疑中，他找到了一个他认为不可能被推翻的起点——"我思故我在"（Cogito, ergo sum）。从这一不可动摇的第一原理出发，笛卡尔试图用理性演绎的方法重建整个人类知识体系。笛卡尔的方法论对欧洲大陆哲学的影响极为深远，它所蕴含的"人的理性具有独立于传统的认识世界的充分能力"这一信念，与启蒙思想对理性权威的强调一脉相承。

这两条求知路径——培根的经验归纳和笛卡尔的理性演绎——虽然在方法上几乎对立，但它们的共同点是拒绝了"经典文本已经够用"的经院传统。世界是可以被重新认识的，不需要依赖任何权威的注释。
}

\subsection{启蒙的思想家与他们的方案}

17世纪中期至18世纪晚期，欧洲知识界经历了一场从"以经文和教会教导为依据判断一切"向"以理性观察和论证为依据判断一切"的深刻转变。这一转变在哲学上被称为"启蒙运动"。

启蒙思想家的共同信念是：通过理性的运用和对自然界的科学研究，人类能够逐步改善自身的命运——包括政治制度的改良。这一信念并不是一种对"进步"的盲目信仰，而是一种基于牛顿物理学成就的总结：既然自然界的运行可以被人认识并加以系统描述，为什么人类社会的运行规则不能被同样地加以分析和改进？

约翰·洛克在1689年出版了《政府论两篇》和《人类理解论》。在政治哲学方面，他系统地论证了政府合法性不是来自神的授予或血缘世袭——而是来自被统治者的同意。私有财产不是国王赐予的，而是早于政府存在的自然权利。如果国王侵犯了这种自然权利——人民有权收回委托给他的权力。美国《独立宣言》起草人杰斐逊在1776年几乎原文照搬了洛克关于"生命、自由和财产"的自然权利表述。

孟德斯鸠在1748年出版的《论法的精神》中分析了欧洲各国不同的政体结构，提出了一项后来在1787年美国制宪会议上被直接采纳的理论预设：立法、行政和司法三种权力应当由三个互不隶属的机构分别掌管，以阻止任何一方的权力膨胀为绝对性的统治权。"权力只有靠权力来制约"——不是依靠统治者的仁慈和德行来保障。

卢梭在1762年出版的《社会契约论》以一句振聋发聩的开场白开始："人生而自由，但无往不在枷锁之中。"卢梭质疑的不是"枷锁"的存在——他认为人在社会中建立制度和政府规范是完全必要的。他质疑的是枷锁所体现的不平等的权力是否具有道德的合法性。他提出政治权力的唯一合法来源是全体人民的共同意志——主权在民。

\subsection{启蒙的物质基础：印刷品、咖啡因和谈话}

启蒙思想的传播并不是因为某些伟大人物写下了一些正确的主张，然后大众默默跟进。文本的传播需要载体和受众，而受众的形成依赖于特定的社会空间。

实验室与研究院：1660年伦敦建立了"增进自然知识的皇家学会"，1666年巴黎创立了"法兰西王家科学院"，柏林和圣彼得堡也相继仿效。这些机构定期举行会议、出版会刊、资助实验并接收来自全欧洲的通信和论文。皇家学会的《哲学汇刊》自1665年开始发行，是最早的科学期刊之一。学会和期刊提供了一种制度化的知识积累机制——一项试验结果不必在研究者死后被弟子发现和传播，而是可以在几个月内跨过英吉利海峡被其他研究者检索、质疑和验证。

印刷机：到18世纪初，印刷术已经运作了两百多年。书籍、期刊、报纸和小册子的生产成本已大幅下降。伏尔泰的著作在当时的印制和发行量使他能够突破首都圈子，其传播范围足以覆盖相当广泛的省域识字人群。狄德罗和达朗贝尔主导的《百科全书》（1751—1772年陆续出版）是印刷业和启蒙运动最有雄心的共同作品——它试图将人类已知的全部知识汇总于一套书中。

载体：印刷机。到18世纪初，德国发明家古登堡的印刷术已经在欧洲运作了两百多年。书籍、期刊、报纸和小册子的生产成本已大幅下降。伏尔泰的著作在当时的印制和发行量使他的政治影响能够突破首都圈子深入到相当广泛的省域识字人群中。

受众：识字率的提升。三十年战争的互相杀戮和宗教冲突使欧洲知识界普遍厌恶为神学教义继续从事流血活动。阅读和讨论不再局限于修道院、大学和宫廷。伦敦的咖啡馆在18世纪初有好几百家，花一个便士买杯咖啡就可以在里面坐一个下午，翻阅当天的报纸并与邻座辩论政治。法语中的"沙龙"在这时已不仅仅指客厅，而是指由开明贵族或富商资助的、定期聚集知识分子进行自由讨论的文化场所。

这些面对面的对话和纸面上的文字交互，构成了启蒙运动最底层的社会基建。不是因为"思想家写得精彩，人们就信服了"——而是因为人们自己聚在咖啡馆和沙龙里，读同一篇文章、讲同一件事、讨论几乎同样的问题，开始对旧制度的正当性产生了一个共同的、自发的怀疑。

\subsection*{A组：内容复述与要点整理}

\begin{enumerate}[label=A\arabic*.]
  \item 哥白尼的日心说对当时欧洲人的世界观构成了什么挑战？
  \item 伽利略为什么在1633年被宗教裁判所审判？他提交给法庭的证据是什么？
  \item 牛顿的物理学对启蒙思想家的政治思维产生了什么间接影响？
  \item 洛克、孟德斯鸠和卢梭各自提出了什么样的政治主张？
  \item 18世纪启蒙思想的传播依赖了哪些社会条件？
\end{enumerate}

\subsection*{B组：深入分析与扩展讨论}

\begin{enumerate}[label=B\arabic*.]
  \item \textbf{概念分析题：}什么是"自然权利"？洛克主张自然权利先于政府存在，政府不得侵犯。这一观念和美国1776年《独立宣言》之间存在什么样的文本和论证关联？这种"自然权利"的观念对中国传统的法律思想构成了什么挑战？
  \item \textbf{比较分析：}孟德斯鸠的三权分立方案和法国同时期君主专制的运作机制之间形成了一组对比——权力分散和权力集中的逻辑各建立在一套什么样的预设之上？在中国古代史中，是否出现过类似的对权力集中的思想性质疑或对立制度设计？
  \item \textbf{论述题}（500-600字）：咖啡馆在18世纪的伦敦被称为"便士大学"。请问当今社会是否存在类似的社会设施，能够以极低的货币门槛将不同阶级、不同知识背景的人群聚集到一个空间中进行自由讨论？如果是——它是什么？它的优点和可能导致的知识质量降低的问题分别是什么？请提供具体的例子。
\end{enumerate}

\vspace{0.5cm}
\begin{center}
\textcolor{blue!70!black}{\large —— 从凡尔登到威斯特伐利亚，从卡诺莎到启蒙运动的文字——欧洲的权力结构史和思想结构史，你已经看见完整的贯穿线了。——}
\end{center}

\end{document}
